% Part: many-valued-logic
% Chapter: syntax-and-semantics
% Section: formulas

\documentclass[../../../include/open-logic-section]{subfiles}

\begin{document}

\olfileid{mvl}{syn}{fml}

\olsection{\usetoken{P}{formula}}

\begin{defn}[الصيغة]
\ollabel{defn:formulas}
مجموعة~$\Frm[L]$، وهي مجموعة~\emph{!!{formula}s} لغة القضايا
$\Lang L$، تعرف بالاستقراء كما يأتي:
\begin{enumerate}
\item كل~!!{propositional variable}~$\Obj p_i$~!!{formula} ذرية.
\item كل رابط ذي~$0$ من المواضع في~$\Lang L$، أي كل ثابت قضوي،
!!{formula} ذرية.
\item إذا كان~$\star$ رابطًا ذا~$n$ من المواضع في~$\Lang L$، وكان
$n>0$، وكانت~$!A_1$، \dots،~$!A_n$~!!{formula}s، كانت
$\star(!A_1, \dots, !A_n)$~!!a{formula}.
\tagitem{limitClause}{ولا يعد~!!a{formula} شيء سوى ما بني بهذه البنود.}{}
\end{enumerate}
فإذا كان~$\star$ ذا موضع واحد، أوجزنا غالبًا~$\star(!A_1)$ إلى
$\star !A_1$. وإذا كان~$\star$ ذا موضعين، أوجزنا
$\star(!A_1,!A_2)$ إلى~$(!A_1 \star !A_2)$.
\end{defn}

وكعادتنا نحذف زوج الأقواس الخارجي حيث لا يقع لبس.

\begin{ex}
  في اللغة القياسية~$\Lang{L_0}$، تعد
  $\Obj p_1 \lif (\Obj p_1 \land \lnot \Obj p_2)$ صيغة.
  وأما في لغة منطق الضرب فتكتب
  $\Obj p_1 \lif (\Obj p_1 \odot \lnot \Obj p_2)$.
  وإذا زدنا في اللغة الرابط الأحادي~$\triangle$، دخلت فيها أيضًا
  صيغ نحو~$\triangle (\Obj p_1 \land \Obj p_2) \lif
  (\triangle \Obj p_1 \land \triangle \Obj p_2)$.
\end{ex}

\end{document}
